\section{Discussion}
\label{sec:discussion}

This study compared a supervised object detection model (YOLOv8) with a multimodal vision language model (LLaVA 1.6) for automated housing inspection using identical datasets, inspection categories, and evaluation metrics. By evaluating both approaches under the same experimental conditions, this research addresses a gap in the literature, where supervised object detectors and multimodal foundation models are typically evaluated separately using different datasets and evaluation protocols. The results demonstrate that YOLOv8 consistently achieved superior performance for inspection category classification, whereas LLaVA demonstrated strengths in reasoning about the overall condition of a property. Together, these findings suggest that the two model paradigms perform best at different levels of the inspection process.

\subsection{Comparison with Previous Research}

The superior category level performance of YOLOv8 is consistent with previous research on building defect detection, infrastructure monitoring, and automated damage assessment, where supervised object detection models achieve strong performance when sufficient annotated training data are available \cite{kang_improvement_2025,lee_automated_2024}. These studies similarly report that explicit supervision enables models to learn domain specific visual characteristics required for accurate localization and classification of structural defects.

The findings also agree with recent research on multimodal foundation models. Previous studies have shown that vision language models possess strong visual reasoning capabilities and demonstrate promising zero shot performance across different domains \cite{jiang_mmad_2025,wen_autonomous_2024,malla_enhancing_2025}. However, these studies also conclude that foundation models become less reliable when applied to specialized inspection tasks requiring precise category discrimination. The results obtained in this study support these observations. Although LLaVA benefited from prompt engineering and correctly interpreted several inspection scenarios, its category level performance remained substantially below that of YOLOv8.

An important contribution of this research is the controlled comparison between both model paradigms. Previous studies generally evaluate supervised object detectors and vision language models independently, making direct comparison difficult. By evaluating both models under identical experimental conditions, this study demonstrates that the observed differences are primarily related to the underlying learning paradigm rather than differences in datasets or evaluation methodology.

\subsection{Interpretation of the Results}

Several factors explain the observed performance differences. YOLOv8 was trained directly on the harmonized inspection categories using bounding box annotations, enabling it to learn visual representations that are specific to housing related defects. In contrast, LLaVA relied entirely on large scale multimodal pretraining and prompt based reasoning without additional task specific learning.

The category specific results further support this interpretation. Categories with distinctive visual characteristics, such as asbestos, achieved relatively high recall, whereas categories such as mold and wear remained considerably more difficult for both models. These categories exhibit greater visual variation and were represented by fewer training examples, making accurate classification inherently more challenging, which is consistent with previous work on automated building defect detection \cite{kang_improvement_2025}.

One of the most important findings of this research is the distinction between category level and property level performance. YOLOv8 consistently achieved the highest Precision, Recall, and F1 scores for individual inspection categories, demonstrating its suitability for detailed defect detection and localization. In contrast, the rule guided LLaVA configuration correctly identified the overall inspection outcome for all evaluated properties despite its lower category level performance. Similar observations have been reported in recent work on multimodal inspection systems and industrial anomaly detection, where vision language models provide stronger semantic reasoning than fine grained localization \cite{jiang_mmad_2025,malla_enhancing_2025}.

These findings indicate that the two approaches should not necessarily be considered competing alternatives. Instead, they appear to operate most effectively at different levels of abstraction within the inspection process. Supervised object detection excels at identifying and localizing individual defects, whereas multimodal foundation models are better suited for interpreting inspection results at the property level, summarizing findings, and supporting higher level decision making.

Prompt engineering improved the consistency of LLaVA and reduced output parsing errors, but the overall improvement remained limited. This observation is consistent with previous research showing that prompt engineering can improve the performance of vision language models but cannot fully compensate for the absence of domain specific supervision \cite{malla_enhancing_2025}.

\subsection{Transferability from Medical Imaging}

One of the motivations for this research was to investigate whether methodological developments from medical imaging could be transferred to housing inspection. Medical imaging has demonstrated that accurate detection of subtle abnormalities depends not only on model architecture but also on standardized evaluation procedures, carefully annotated datasets, and appropriate image resolution \cite{litjens_deep_2017,kamraoui_longitudinal_2022,rana_yolo_2025}.

Several of these principles were incorporated into this study. Public datasets were harmonized into a unified inspection framework, standardized evaluation metrics were applied throughout the experiments, and the input resolution was increased after qualitative analysis revealed that smaller defects were frequently missed. These methodological decisions improved the consistency of the evaluation and contributed to a more reliable comparison between both model paradigms.

Although housing inspection and medical imaging differ substantially as application domains, both require the identification of relatively small visual abnormalities under variable imaging conditions. The findings therefore suggest that transferring methodological practices from mature computer vision domains can improve the development and evaluation of automated housing inspection systems.

\subsection{Limitations}

Several limitations should be considered when interpreting the results.

First, the harmonized dataset was constructed by combining multiple publicly available datasets originating from different domains. Although careful label mapping was performed, differences in image composition, annotation quality, and acquisition procedures may still have influenced model performance.

Second, class imbalance remained an important challenge throughout the experiments. Categories such as wear contained considerably fewer examples than damage and crack, which likely contributed to the lower recall observed for these categories. Future research would benefit from larger and more balanced housing inspection datasets.

Third, this study evaluated only one supervised detector and one multimodal foundation model. Although both models were selected because they represent widely adopted approaches within their respective paradigms, the findings cannot automatically be generalized to newer detector architectures or larger vision language models.

Finally, the practical evaluation was performed using inspection images from a single housing corporation. While this strengthens the ecological validity of the experiments, additional validation using data from multiple organizations would improve the generalizability of the findings.

\subsection{Practical Implications}

The findings have several practical implications for housing inspection.

For detailed defect detection and localization, supervised object detection currently represents the most reliable solution because it combines accurate classification with precise localization of inspection relevant defects, consistent with previous studies in automated building inspection \cite{kang_improvement_2025}. This makes YOLOv8 well suited for supporting inspectors during property assessments and documenting the location of detected defects.

However, the results also demonstrate that vision language models provide complementary capabilities. Whereas YOLOv8 performed best at the inspection category level, LLaVA demonstrated stronger performance at the property level by correctly identifying the overall inspection outcome after prompt refinement. This indicates that the strengths of the two paradigms lie at different stages of the inspection workflow.

Rather than replacing supervised object detection, vision language models should therefore be viewed as complementary technologies. A promising practical direction is the development of hybrid inspection systems that combine the strengths of both approaches. In such a workflow, YOLOv8 could first detect and localize individual defects using bounding boxes, after which a vision language model could interpret these detections, assess the overall condition of the property, generate inspection summaries, and support inspectors during decision making. By combining accurate visual localization with higher level semantic reasoning, such systems have the potential to provide more comprehensive decision support than either model can achieve independently.

This hybrid approach represents the logical continuation of the present research and forms the basis for future work, in which the combined system will be evaluated in real world housing inspection scenarios to determine whether it improves inspection accuracy, efficiency, and practical usability.

